\section{A signed horizon and an integral rescaling}
\label{sec:signed}

\subsection{The exact arithmetic interface}

We retain the 210-variable parent's literal signed Pell, binomial, and
Chinese remainder construction~\cite{growth210}. Its horizon, originally
\(t=16+4z^2\), is changed to
\begin{equation}\label{eq:horizon}
 t=2+4z,\qquad z\in\Z.
\end{equation}
Write \(d=2^{19}\). The three standalone growth equations remain
\begin{equation}\label{eq:growthrows}
 Z=tF,\qquad X^2-16Z^2-8ZF=1,\qquad F=dt(2u+1).
\end{equation}
At this stage \(t,F,X,Z,u\) are signed. Every displayed equation has degree
at most two. The Pell equation obtained by substituting the first row is
\(X^2-((4t+1)^2-1)F^2=1\); this substituted expression is for the proof,
not an additional higher-degree residual in the quartic.

The parent also retains an even anchor \(C\), a binomial selector \(B\),
and the identities
\begin{equation}\label{eq:radix}
 B\eta=C+F+1,\qquad D=2(C+F+1).
\end{equation}
Its ordinary polynomial power certificates have the semantic consequences
\begin{equation}\label{eq:power-interface}
 q=D^x,\qquad q^t,\qquad q^{y+1},\qquad
 (q+1)^{F+t},\qquad (q+1)^C,
\end{equation}
as well as the other requested powers in the exact source. A notation
such as \(q^t\) in this interface means that the retained polynomial
certificate has a proposed integral output for this base and exponent.
It is not a variable-exponent expression in the final polynomial.
For \(\abs q\ge2\), the certificate forces its exponent to be nonnegative
and its output to be the true power.

The beta endpoints are
\begin{equation}\label{eq:beta-endpoints}
 \begin{aligned}
 y&=F(2v+1),\\
 x&=F+1+(y+1)r_0,\\
 x&=1+\bigl(1+(t+1)y\bigr)r_t.
 \end{aligned}
\end{equation}
The letters in this interface are semantic names for fixed source
coordinates or polynomial abbreviations. The final variable list is
given independently in the literal equation data.

\subsection{Recover the signs before applying a sharp growth bound}

\begin{lemma}[Preliminary bound]\label{lem:preliminary}
Every integer solution of \eqref{eq:horizon} and the last row of
\eqref{eq:growthrows} satisfies
\[
 t\ne0,-1,\quad \abs t\ge2,\quad F\ne0,\quad
 2^{20}\mid F,\quad \abs F\ge d\abs t>\abs t.
\]
\end{lemma}
\begin{proof}
The progression in \eqref{eq:horizon} is exactly the integers congruent to
2 modulo 4. Thus \(\vTwo(t)=1\). The integer \(2u+1\) is nonzero and
odd, so \(\vTwo(F)=20\), and its absolute value is at least one.
\end{proof}

The parent's square frame includes the actual growth coordinate \(F\).
Consequently its parameter \(W\) is at least \(2+F^2>258\), even though
\(t^2\) is now permitted to be smaller than 256. This supplies the
original common-frame threshold before signs are recovered. The cap and
orientation proof uses squares and absolute values: it classifies the
anchor and all principal Pell pairs at their canonical signed indices
without assuming \(t>0\) or \(F>0\). The new standalone pair in
\eqref{eq:growthrows} is independent of this common frame and needs no
new frame atom. All actual bases, nominal indices, and outputs retained
from the parent remain in the frame.

\begin{lemma}[Noncircular sign recovery]\label{lem:signs}
With the retained canonical-index and exact-power interface, every zero
of the changed source has
\[
 t\ge2,\quad x\ge2,\quad y\ge F>0,\quad C>0,
 \quad D>2C>0,\quad q>0.
\]
\end{lemma}
\begin{proof}
Both \(C\) and \(F\) are even. Thus \(C+F+1\) is odd and nonzero;
\eqref{eq:radix} makes \(B,\eta\) odd and nonzero, and
\(\abs D\ge2\). The request \(q=D^x\) gives \(x\ge0\).

Set \(g=1+(t+1)y\). By Lemma~\ref{lem:preliminary} and
\eqref{eq:beta-endpoints},
\[
 \abs g\ge\abs{t+1}\abs y-1
       \ge\abs F-1>1,
\]
because \(t\ne-1\). If \(x=0\), the last endpoint would give
\(gr_t=-1\), a contradiction. Therefore \(x\ge1\), and
\(\abs q\ge2\). The retained request with exponent \(t\), exactly
\(t\) rather than \(t+1\), now forces \(t\ge0\). The progression
excludes zero, so \(t\ge2\).

The request with exponent \(y+1\) gives \(y+1\ge0\). A negative \(y\)
would have \(y\le-\abs F<-1\). Hence \(y\ge\abs F>0\). If \(x=1\),
the first endpoint says \((y+1)r_0=-F\), impossible since
\(0<\abs F\le y<y+1\). Thus \(x\ge2\), \(\abs q\ge4\), and
\(\abs{q+1}\ge3\).

The requested \((q+1)\)-powers now imply \(F+t\ge0\), \(C\ge0\),
and nonnegativity of every other actual CRT index. A negative \(F\)
would give \(F+t=t-\abs F<0\), so \(F>0\). The anchor proof already
excludes \(C=0\), yielding \(C>0\). Formula~\eqref{eq:radix} gives
\(D>2C>0\) and then \(q=D^x>0\).
\end{proof}

This order matters. On a negative branch the positive Pell center would
be \(\abs{4t+1}=4\abs t-1\). One cannot apply a lower bound proved for
\(4\abs t+1\) there. The elementary inequality
\(\abs F\ge d\abs t\) supplies all the initial sign information we need.

\subsection{The sharp positive-center estimate works at the smallest horizon}

For \(a\ge2\), write \(X_n(a),Y_n(a)\) for the ordinary Pell sequences
\[
 X_n(a)+Y_n(a)\sqrt{a^2-1}
       =\bigl(a+\sqrt{a^2-1}\bigr)^n.
\]
The recurrence, rather than the radical notation, gives their integer
definition. It implies
\begin{equation}\label{eq:pellfacts}
 \begin{gathered}
 Y_n(a)\equiv n\pmod{a-1},\qquad
 \vTwo(Y_n(a))=\vTwo(n)\quad(a\text{ odd},\ n>0),\\
 Y_n(a)\ge(2a-1)^{n-1}\quad(n\ge1).
 \end{gathered}
\end{equation}
The valuation identity follows by repeatedly using
\(Y_{2m}=2X_mY_m\), since every \(X_m\) is odd and \(Y_m\) is odd
exactly for odd \(m\). The last inequality follows from
\(Y_{n+1}=2aY_n-Y_{n-1}\) and monotonicity.

By Lemma~\ref{lem:signs}, \(a=4t+1\ge9\). Classification of integral
Pell solutions gives \(F=Y_n(a)\) for a positive \(n\). Since
\(4t\mid F\), the congruence in \eqref{eq:pellfacts} yields \(4t\mid n\).
The exact valuation of \(F=dt(2u+1)\) yields
\(\vTwo(n)=19+\vTwo(t)\). Therefore
\begin{equation}\label{eq:indexbound}
 n\ge dt.
\end{equation}

The unchanged square-shell constructors are
\[
 \begin{aligned}
 I(a,b)&=2((a+b)^2+a),& M(a,b)&=I(a,b)+1,\\
 A(a,b)&=2((a+b)^2+3a+4b+3),&E(a,b)&=A(a,b)+1.
 \end{aligned}
\]
For arguments at most \(m\), each output plus one is at most
\(8(m+1)^2\). All local parameters, and the constants 0 and 1, are
strictly below \(t\) when \(t\ge2\). The source's authenticated maximum
constructor depth is 19. Induction therefore bounds every local label
strictly below
\begin{equation}\label{eq:labelbound}
 L=8^{d-1}t^d,\qquad d=2^{19}.
\end{equation}
Combining \eqref{eq:pellfacts}--\eqref{eq:labelbound},
\begin{equation}\label{eq:sharpbound}
 F\ge(8t+1)^{n-1}\ge(8t)^{dt-1}\ge L^t.
\end{equation}
The last ratio is \(8^{t-1}/t\ge1\), including \(t=2\).

Two older numerical uses of \(t\ge16\) deserve explicit replacement.
For \(\Lambda_t=\operatorname{lcm}(1,\ldots,t)\),
\[
 \Lambda_t\le t!\le t^t,\qquad
 F\ge L^t\ge t^{2t}\ge\Lambda_t^2.
\]
Thus the actual hypothesis of the factorial-free binomial CRT theorem
still holds. Also
\(F\ge L^2\ge2^{8d-6}>2^{32}\), so all coarse lower bounds on the
later radix and beta code remain available. Finally the positive endpoint
gives \(x\ge2+(t+1)y\). Since \(y\ge F\),
\[
 2+t+y+F<x,
\]
already for \(t\ge2\), because \((t-1)y>t\). This validates the
parent's local bound \(2+t+x+y+F<2x\).

With these replacements, the finite-product separation theorem uses
precisely \(F\ge L^t\); the binomial decoder uses \(D>2C\), the
canonical common frame, and the retained remainder bound; and the beta
and Horn-flow arguments recover a finite \(\Ccal\) refutation. The
logical calibration then supplies a finite syntactic ZFC refutation.

\subsection{Completeness on a cofinal progression}

Starting from a finite calibrated refutation, choose a sufficiently large
positive horizon
\begin{equation}\label{eq:cofinalt}
 t\equiv2\pmod{2^{23}}.
\end{equation}
This is a cofinal subprogression of \eqref{eq:horizon}; it has
\(\vTwo(t)=1\), and binary digit 20 of \(t\) is zero. All finite proof
length and parameter bounds can be met.

Put \(n_0=dt\), \(y_0=Y_{n_0}(4t+1)\),
\(a_0=X_{n_0}(4t+1)\), and \(b_0=y_0/(dt)\).
The preceding congruence and valuation facts show that \(b_0\) is an
odd integer. Since \(n_0\) is even, \(a_0\equiv1\pmod4\).

For every integer \(a_0\equiv1\pmod4\), the map
\(s\mapsto Y_s(a_0)\) permutes the residue classes modulo each power
of two. Here is the needed induction: for \(h=2^r\), the norm and
valuation identities give
\[
 X_h(a_0)\equiv1\pmod{2^{2r+2}},\qquad
 Y_h(a_0)\equiv2^r\pmod{2^{r+1}}.
\]
The addition formula yields
\(Y_{s+2^r}(a_0)-Y_s(a_0)\equiv2^r\pmod{2^{r+1}}\).
Thus the two lifts of every index class produce the two different lifts
of its image. Starting modulo two proves the assertion, together with
periodicity and arbitrarily large indices in a prescribed class.

Choose \(\ell>23\) with \(2^\ell>t\). The two odd integers \(t/2\)
and \(b_0\) are invertible modulo \(2^{\ell-20}\). Choose an arbitrarily
large positive odd \(s\) satisfying
\[
 (t/2)b_0Y_s(a_0)\equiv1\pmod{2^{\ell-20}}.
\]
Then set
\[
 F=Y_{n_0s}(4t+1)=y_0Y_s(a_0),\quad
 X=X_{n_0s}(4t+1),\quad Z=tF,\quad
 u=\frac{F/(dt)-1}{2}.
\]
All witnesses are integral, \eqref{eq:growthrows} holds, and
\(F\equiv2^{20}\pmod{2^\ell}\). The binary supports of \(F\) and \(t\)
are disjoint. The elementary binary-carry criterion consequently gives
\(B=\binom{F+t}{t}\) odd.

Keep the beta multiplier independent. Choose \(2v+1\) to be a positive
odd multiple of every odd prime at most \(t\), and sufficiently large
that \(y=F(2v+1)\) exceeds all finitely many proof-state values. Every
prime at most \(t\) divides \(y\). If a prime divided both \(1+iy\)
and \(1+jy\), \(1\le i<j\le t+1\), it would divide \(j-i\) and be
coprime to \(y\), a contradiction. Hence the beta moduli are pairwise
coprime. Full factorial divisibility of \(y\) is unnecessary.

The witness order is now acyclic: choose beta code \(x\), then \(B\) and
\(N=B^x\), the nonanchor CRT representatives, an admissible even anchor
\(C\) larger than all other nominal indices, the affine radix \(D\),
the true power outputs, and finally the common Pell frame and cap.
The retained CRT theorem applies by \eqref{eq:sharpbound}. This constructs
a zero from every finite refutation and proves completeness of the changed
source.

\subsection{From the semantic source to the two literal quartics}

The preceding horizon change is an existence-preserving semantic change.
It is not asserted to be a bijection with the old square slice.
After that change, all remaining transformations are explicit bijections
on integer zero sets.

In the parent ordering, replace row 125 by \(x_0-2-4x_{178}=0\).
Delete that row, use the new coordinate \(\xi_0=x_{178}\), and substitute
\(x_0=2+4\xi_0\). Keep the other coordinate order, deleting old coordinate
178. For the sparse profile additionally substitute \(F=64\xi_2\) and
divide the last growth-divisibility residual by 64. Lemma~\ref{lem:preliminary}
guarantees the inverse \(\xi_2=F/64\) is integral at every zero.
All residuals remain quadratic. The final growth residual is
\begin{equation}\label{eq:lastgrowth}
 \xi_2-16384-32768\xi_0-32768\xi_{208}
                   -65536\xi_0\xi_{208}.
\end{equation}
Its square contributes coefficient \(2^{32}\) at
\(\xi_0^2\xi_{208}^2\). The last coordinate occurs in no other residual,
so this coefficient cannot cancel. Exact collection verifies that no
coefficient is larger and that precisely 5,486 monomials remain.

For the lower-height companion use \(F=128\xi_2\), whose inverse is
likewise integral. A retained residual with coefficient \(16384\) at
\(\xi_2^2\) is used to remove the corresponding square terms from two
other rows by integer row additions. With zero-based final row indices, the operations are
\[
 R_{98}\leftarrow R_{98}+4R_{135},\qquad
 R_{124}\leftarrow R_{124}+6R_{135},
\]
with \(R_{135}\) retained. The certificate lists every coefficient. Such a triangular residual-basis operation has an integer
inverse; it changes the sum-of-squares polynomial while preserving its
zero set. The resulting polynomial has 6,654 monomials and height
\(2^{30}\). Both statements are established by full coefficient
collection, not by estimates of input literals.
